% Complete conversion of audited historical support report; no KMS theorem imports.
\section{Historical support branches, ideal completion, and measured observations}\label{historical-support-historical-support-branches-ideal-completion-and-measured-observations}

This chapter gives the complete branch, ideal-completion, atomic-kernel,
and operator-algebra proofs extracted and continued from the historical
Split-Zero support program. Its mathematical report and independent audit
are, respectively, \path{logbook/historical_next_responses_20260912.md} and
\path{logbook/frame_branch_independent_audit_20260912.md}.
The equation identifiers B, I, M, A and O are local to this chapter.

\textbf{Source locators and comparison.} Exact complete response bodies are staged at
\path{sources/historical/next_responses/preboolean_prime_branches.md}
(source lines 23137--23956; response heading 23104; exported timestamp
7/6/2026, 12:23:35 PM),
\path{sources/historical/next_responses/preboolean_arithmetic_branch_maps.md}
(source lines 25162--25715; heading 25058; timestamp 7/6/2026, 7:58:41 PM), and
\path{sources/historical/next_responses/pointfree_frame_arithmetic_claim.md}
(source lines 25754--26221; heading 25721; timestamp 7/6/2026, 10:01:27 PM).
Their shared original source has SHA-256
\nolinkurl{8d1780a24df3bb31db981b5e437558de31e6d94bf21bb54103a0a336d6a73091}.
The respective excerpt SHA-256 values are
\nolinkurl{f5f11dd6b26a3a9d6a4c77cf2f7b86898759a19fa740c993d5749c8d47c02618},
\nolinkurl{e84873955876b2be5f7738f746b4e133fc479cfe726bcb85d477f9aec7d6d10e}, and
\nolinkurl{57c6662732b723e451a495f0224507f61e3a969956178018af7633ae06797ac0}.
For each excerpt, its first line is the first indicated original source
line, so every original locator has the exact excerpt locator obtained by
subtracting that first source line and adding one.

The core corpus comparison is
\path{split_zero_projective_monads_surcomplex/split_support_geometry_arithmetic_curve_v11.tex},
SHA-256 \nolinkurl{5fa6a55aa9d1a2e57599e206087d3d2466b25dd99fa1a90d5f86367aafab1b4d}.
Its lines 1772--2054 already prove the ordinary prime-filter branches and
their arithmetic maps; lines 2001--2018 explicitly impose arbitrary-join
preservation for locale points. The support and measured realizations were
also compared at lines 5006--5096 and 5346--5378. Accordingly the ordinary
branch classification is a surviving historical construction, while the
ideal quotient, explicit descent witnesses, atomic kernel and full
\(L^\infty\) correspondence below supply the additional proved maps.

The final historical response, source lines 25982--26020 (excerpt lines
229--267), equates ordinary amplitude-fixed semiring branches with frame
points without imposing arbitrary-join preservation. The exact correction
is the classification (B1)--(B5), the frame quotient (I2), and the
point-descent criterion (I3)--(I10). The measure-algebra construction then
exhibits the kernel pairs that ordinary branches can distinguish. No
external KMS classification or Morita equivalence is used as a theorem
of this chapter.

\subsection{Exact ordinary and continuous Split-Zero branches}\label{historical-support-exact-ordinary-and-continuous-split-zero-branches}

The results below are stated in classical set theory. The ultrafilter construction explicitly uses Zorn's lemma. No claim about existence or nonexistence of points in a constructive foundation is inferred from that use of choice; such an internal comparison would also have to specify its truth-value object and its available joins.

Let \(R\) be a nonzero commutative ring and let \(L\) be a nontrivial bounded distributive lattice. Retain the source carrier and all its exceptional coordinates:

\[
G_L(R)=\{(0,\lambda):\lambda\in L\}\cup\{(r,1_L):r\in R\}\subseteq R\times L,
\]

\[
(r,\lambda)+(s,\mu)=(r+s,\lambda\vee\mu),\qquad
(r,\lambda)(s,\mu)=(rs,\lambda\wedge\mu),
\]

\[
\tau_L=(0,0_L),\qquad e_L=(0,1_L),\qquad1=(1_R,1_L).
\]

Write \(G(R)=\{\tau\}\sqcup R^\bullet\), with supported zero \(e=0_R^\bullet\) distinct from global zero \(\tau\). The indicated subset is the carrier throughout.

Here \(\mathbb B=\{0,1\}\) has join and meet as its semiring operations.
The supported copy in \(G(R)\) has
\(r^\bullet+s^\bullet=(r+s)^\bullet\) and
\(r^\bullet s^\bullet=(rs)^\bullet\); its global zero satisfies
\(\tau+x=x+\tau=x\) and \(\tau x=x\tau=\tau\).
The map \(G_{\mathbb B}(R)\to G(R)\) defined by
\((0,0)\mapsto\tau\), \((r,1)\mapsto r^\bullet\) is a bijection and
preserves these operations in each support case, hence is a unital semiring
isomorphism. To verify the general carrier, the displayed operations are
the restrictions of the product semiring \(R\times L\). This subset is
closed: a sum of nonzero amplitude has top support because at least one
input has nonzero amplitude; a product of nonzero amplitude has two
nonzero input amplitudes and hence top support. Every zero-amplitude
output belongs to the zero fibre. The zero and unit belong to the subset.
Thus associativity, commutativity and distributivity restrict from the
product, including all supported cancellations. All homomorphisms below
preserve zero and unit. Write \(z_\lambda=(0,\lambda)\).

\textbf{Theorem B1 --- all ordinary amplitude-fixed branches.} There is a natural bijection

\[
\operatorname{Hom}_{\mathrm{DLat}_{0,1}}(L,\mathbb B)
\xrightarrow{\sim}
\{\beta:G_L(R)\to G(R)\text{ unital semiring map}:\beta(r,1_L)=r^\bullet\}.
\tag{B1}
\]

It sends \(\epsilon\) to

\[
\beta_\epsilon(r,1_L)=r^\bullet,\qquad
\beta_\epsilon(0,\lambda)=
\begin{cases}\tau&\epsilon(\lambda)=0,\\ e&\epsilon(\lambda)=1.\end{cases}
\tag{B2}
\]

The map is surjective as a semiring map; its exact kernel congruence is equality of amplitudes together with equality of selected support values:

\[
(r,\lambda)\sim_\epsilon(s,\mu)
\iff r=s\ \text{and}\ \epsilon(\lambda)=\epsilon(\mu).
\tag{B3}
\]

\textbf{Proof.} The source formulas agree at \((0,1_L)\) since \(\epsilon (1_L)=1\). They preserve zero and unit. On the zero fibre, preservation of addition and multiplication is precisely preservation of join and meet. If one argument has nonzero amplitude, its lattice coordinate is \(1_L\); adding a zero-fibre argument leaves it unchanged and multiplying by that argument gives \((0,\lambda )\), matching respectively addition by \(\tau\) or \(e\), and multiplication by \(\tau\) or \(e\) in \(G(R)\). With both arguments at full support, even if their sum or product cancels to amplitude zero, both outputs are supported and ring identities prove preservation. This verifies all exceptional products and cancellations.

Conversely, \((0,\lambda )\) is additively idempotent. The only additive idempotents of \(G(R)\) are \(\tau\) and \(e\): a supported \(r^\bullet+r^\bullet=r^\bullet\) implies \(2r=r\) in the ring, hence \(r=0\). Therefore every amplitude-fixed \(\beta\) sends the zero fibre into \(\{\tau,e\}\). Identifying this two-element subsemiring with \(\mathbb B\) gives a map \(\epsilon\) preserving bottom, top, meet and join. Formula (B2) reconstructs \(\beta\); both constructions are inverse. Every amplitude and \(\tau\) occurs in the image, and inspecting (B2) gives exactly (B3). \(\square\)

Now let \(L\) be a frame: all joins exist and finite meets distribute over arbitrary joins. Its zero fibre has specified infinitary joins

\[
\bigvee_{i\in I}^{\rm fib}(0,\lambda_i)=(0,\bigvee_{i\in I}\lambda_i).
\]

\textbf{Theorem B2 --- exactly the continuous branches.} Under (B1), the subset of branches preserving all these zero-fibre joins is exactly

\[
\operatorname{Hom}_{\mathrm{Frm}}(L,\mathbb B).
\tag{B4}
\]

Indeed the corresponding equality for (B2) says exactly \(\epsilon(\bigvee_i\lambda_i)=\bigvee_i\epsilon(\lambda_i)\), and finite meets are already preserved. This is the precise additional structure required by the transcript's final frame formulation. Its ordinary algebraic maps retain a larger, explicitly identified category.

The classification extends to every map between original lattice-split carriers. For commutative unital rings \(R,A\) and nontrivial bounded distributive lattices \(L,M\),

\[
\operatorname{Hom}_{\mathrm{Semi}}(G_L(R),G_M(A))
\cong\operatorname{Hom}_{\mathrm{Ring}}(R,A)
\times\operatorname{Hom}_{\mathrm{DLat}_{0,1}}(L,M).
\tag{B5}
\]

The pair \((h,f)\) gives \((r,\lambda )\mapsto (h(r),f(\lambda ))\). Conversely an arbitrary semiring map sends each additively idempotent \(z_\lambda\) to amplitude zero, say \(z_{f(\lambda)}\); this preserves finite joins, meets and bottom. The identity \(e_L=(1_R,1_L)+(-1_R,1_L)\) makes its image have top support, because one summand is the target unit; additive idempotence makes its amplitude zero. Thus \(e_L\mapsto e_M\) and \(f\) preserves top. Now \((r,1_L)+e_L=(r,1_L)\) forces its image to have top support. Write that image \((h(r),1_M)\). The supported operations give a unital ring map \(h\), and the formulas recover every source element, consistently at \(e_L\). Coordinate composition proves naturality and both inverse constructions. In particular the entire two-factor datum is recoverable from the semiring map; scalar amplitude and support are not freely interchanged.

\subsection{The exact frame comparison retaining all ordinary branches}\label{historical-support-the-exact-frame-comparison-retaining-all-ordinary-branches}

For the underlying bounded distributive lattice of a frame \(L\), let \(\operatorname{Idl}(L)\) be the set of lattice ideals. An ideal is a nonempty lower set closed under finite joins. Its bottom is \(\{0_L\}\), its top is \(L\), its finite meet is intersection, and the join of a family of ideals is

\[
\bigvee_\alpha I_\alpha=
\{x\in L:x\le y_1\vee\cdots\vee y_n
\text{ for some finite }n\ge0,\ y_j\in\bigcup_\alpha I_\alpha\}.
\tag{I1}
\]

The empty finite join is \(0_L\); this retains the bottom ideal in every formula.

\textbf{Theorem I1 --- frame quotient and principal-ideal section.} The maps

\[
q:\operatorname{Idl}(L)\to L,\quad I\mapsto\bigvee I,
\qquad\eta:L\to\operatorname{Idl}(L),\quad\lambda\mapsto\downarrow\lambda
\tag{I2}
\]

have the following exact properties:

\begin{enumerate}
\def\labelenumi{\arabic{enumi}.}
\tightlist
\item
  \(\operatorname{Idl}(L)\) is a frame and \(q\) is a surjective frame homomorphism.
\item
  \(\eta\) is an injective bounded-lattice homomorphism, \(q\eta=\mathrm{id}_L\), and
  \(q(I)\le\lambda\ \Longleftrightarrow\ I\subseteq\eta(\lambda)\).
\item
  The kernel congruence of \(q\) is exactly \(I\sim J\ \Longleftrightarrow\ \bigvee I=\bigvee J\). No claim that \(\eta\) preserves arbitrary joins is made.
\end{enumerate}

\textbf{Proof.} Formula (I1) is an ideal and is the least ideal containing every \(I_\alpha\); intersections give infima. To prove frame distributivity, let \(x\in I\cap\bigvee_\alpha J_\alpha\). Choose \(x\le y_1\vee\cdots\vee y_n\) with \(y_j\in J_{\alpha _j}\). In the distributive lattice,

\[
x=x\wedge(y_1\vee\cdots\vee y_n)
=(x\wedge y_1)\vee\cdots\vee(x\wedge y_n).
\]

Each \(x\wedge y_j\) belongs to \(I\cap J_{\alpha _j}\), proving \(I\cap\bigvee_\alpha J_\alpha\subseteq\bigvee_\alpha(I\cap J_\alpha)\); the reverse inclusion follows from containment. Thus \(\operatorname{Idl}(L)\) is a frame.

The supremum of the right side of (I1) equals \(\bigvee_\alpha\bigvee I_\alpha\), proving arbitrary-join preservation by \(q\), including the empty family. Since \(L\) is a frame,

\[
(\bigvee I)\wedge(\bigvee J)
=\bigvee_{i\in I,j\in J}(i\wedge j)
=\bigvee(I\cap J).
\]

The last equality holds because each \(i\wedge j\) is in the intersection, and each \(x\in I\cap J\) is the term \(x\wedge x\). Hence \(q\) preserves finite meets and the top. Principal ideals satisfy \(\downarrow \lambda \cap \downarrow \mu =\downarrow (\lambda \wedge \mu )\) and \(\downarrow \lambda \vee \downarrow \mu =\downarrow (\lambda \vee \mu )\); they also preserve bottom/top. Their suprema are the original elements, proving \(q\eta=\mathrm{id}_L\), injectivity and surjectivity. The adjunction equivalence and the kernel description follow immediately from the definition of supremum. \(\square\)

\textbf{Theorem I2 --- all algebraic branches become frame points before the quotient.} There are mutually inverse bijections

\[
\operatorname{Hom}_{\mathrm{DLat}_{0,1}}(L,\mathbb B)
\cong\operatorname{Hom}_{\mathrm{Frm}}(\operatorname{Idl}(L),\mathbb B),
\tag{I3}
\]

given by

\[
\widehat\epsilon(I)=\bigvee_{\lambda\in I}\epsilon(\lambda)
\quad\text{and}\quad\epsilon=\widehat\epsilon\eta.
\tag{I4}
\]

A frame point \(\widehat\epsilon\) in (I3) factors through \(q\) if and only if its original \(\epsilon\) is a frame point of \(L\). Thus the exact inclusion of point sets is

\[
q^*:\operatorname{Pt}(L)\hookrightarrow\operatorname{Pt}(\operatorname{Idl}(L)),
\quad\epsilon\mapsto\epsilon q,
\tag{I5}
\]

and its image is the points constant on the kernel congruence in Theorem I1.

\textbf{Proof.} In (I4), the value is \(1\) precisely when the ideal meets the prime filter \(\epsilon^{-1}(1)\). For a join of ideals, a value \(1\) is witnessed by an element below a finite join of source elements; finite primeness then gives an element of value \(1\) in one of those source ideals. The reverse implication follows by containment, proving arbitrary-join preservation. For intersections, if \(I\) and \(J\) each have a value \(1\) element \(i,j\), their meet \(i\wedge j\in I\cap J\) has value \(1\); the reverse implication is immediate. Top and bottom follow from \(\epsilon (1)=1\), \(\epsilon (0)=0\). Therefore \(\widehat\epsilon\) is a frame homomorphism.

Every ideal is the join in \(\operatorname{Idl}(L)\) of its principal ideals, so an arbitrary frame point is completely determined by its restriction along \(\eta\), giving both inverse identities. If \(\epsilon\) preserves arbitrary joins, \(\widehat\epsilon(I)=\epsilon(\bigvee I)=(\epsilon q)(I)\). Conversely if \(\widehat\epsilon=\chi q\) for a frame map \(\chi\), then \(\epsilon=\widehat\epsilon\eta=\chi q\eta=\chi\), so \(\epsilon\) is a frame map. A set map constant on the fibres of surjective \(q\) factors uniquely; if that map is a frame map, the induced factor preserves arbitrary joins and finite meets by lifting families through \(q\). This proves the final image claim. \(\square\)

\textbf{Theorem I3 --- original Split-Zero comparison and exact congruence.} Every bounded-lattice map \(f:L\to M\) induces

\[
G_f(R):G_L(R)\to G_M(R),\quad (r,\lambda)\mapsto(r,f(\lambda)).
\tag{I6}
\]

For the nontrivial frames in Theorem I1 this gives a split pair of unital semiring maps

\[
G_L(R)\xrightarrow{G_\eta(R)}G_{\operatorname{Idl}(L)}(R)
\xrightarrow{G_q(R)}G_L(R),\qquad G_qG_\eta=\mathrm{id}.
\tag{I7}
\]

They preserve \(\tau\) and supported \(e\) as distinct elements. The exact kernel of \(G_q(R)\) consists of equal amplitudes and equal suprema of zero-fibre ideals:

\[
(r,I)\sim(s,J)\iff r=s\text{ and }\bigvee I=\bigvee J.
\tag{I8}
\]

\textbf{Proof.} A nonzero amplitude in \(G_L(R)\) has full support and maps to full support because \(f(1_L)=1_M\); hence (I6) lands in the exact target carrier. Finite join and meet preservation establish the semiring identities coordinate by coordinate, including zero-amplitude cancellation. Formula (I7) follows from \(q\eta=\mathrm{id}_L\), and (I8) is equality of the two coordinates of the image. The bottom and top of the nontrivial source and target lattices differ; their images under these bounded maps are again bottom and top, proving the exact zero statements. \(\square\)

Thus an ordinary branch which fails zero-fibre join continuity is a fully explicit frame point \textbf{upstairs}, and the arithmetic map fails to descend through a proved frame quotient \textbf{downstairs}. The original support object and all its algebraic branches remain related by the maps (I2)--(I8).

The two available lifts have an exact comparison, which prevents a split algebra map from being mistaken for a split frame map. Given an ordinary \(\epsilon :L\to \mathbb B\), both

\[
\epsilon q:\operatorname{Idl}(L)\to\mathbb B,
\qquad\widehat\epsilon:\operatorname{Idl}(L)\to\mathbb B
\tag{I9}
\]

are bounded-lattice maps and agree on every principal ideal. The first descends through \(q\) by definition; the second preserves arbitrary joins by Theorem I2. They coincide on all ideals exactly when \(\epsilon\) preserves arbitrary joins. The forward implication follows because \(\epsilon q=\widehat\epsilon\) and surjectivity of the frame map \(q\) force its factor \(\epsilon\) to preserve joins; the converse is (I4). This also yields two amplitude-fixed Split-Zero branch maps out of \(G_{\operatorname{Idl}(L)}(R)\), with exactly the same criterion for equality.

There is a further exact localic structure. The map \(j=\eta q\) is a nucleus on \(\operatorname{Idl}(L)\): it is monotone, \(I\subseteq j(I)\) because each \(i\in I\) lies below \(\bigvee I\), it satisfies \(j^2=j\) because \(q\eta=\mathrm{id}_L\), and it preserves finite meets because both \(q\) and \(\eta\) do. Its fixed ideals are exactly the principal ideals. Their inherited meet and nucleus-applied join

\[
\bigvee^{\mathrm{fix}}_\alpha\downarrow\lambda_\alpha
=j\left(\bigvee^{\operatorname{Idl}(L)}_\alpha\downarrow\lambda_\alpha\right)
=\downarrow\left(\bigvee^L_\alpha\lambda_\alpha\right)
\tag{I10}
\]

identify the fixed-point frame exactly with the original \(L\). This proves the quotient/sub-locale correspondence at the level of operations, while retaining the precise failure of \(\eta\) to preserve the ambient ideal-frame joins.

\subsection{A fully explicit nontrivial frame with algebraic branches and no frame points}\label{historical-support-a-fully-explicit-nontrivial-frame-with-algebraic-branches-and-no-frame-points}

Let \(L\) be the Lebesgue probability measure algebra of \([0,1]\): measurable subsets modulo null symmetric difference. Write \([A]\) for the class of \(A\), keep the original measure \(\mu (A)\), and use union, intersection and complement for finite Boolean operations.
These operations are well defined modulo null sets, since changing their
inputs changes any output only within the union of those null differences.
The exact order is \([A]\le[C]\) if and only if
\(\mu(A\setminus C)=0\).

\textbf{Lemma M1 --- completeness and atomlessness.} This \(L\) is a complete Boolean algebra, hence a frame, and has no atoms.

\textbf{Proof.} For an arbitrary family \(([A_i])_{i\in I}\), let

\[
b=\sup\{\mu(A_{i_1}\cup\cdots\cup A_{i_n}):n\ge0\}.
\]

Choose finite unions \(F_n\) with measures tending to \(b\), and put \(B=\bigcup_n F_n\). Replacing \(F_n\) by the union of the first \(n\) such finite unions makes them increasing without changing their defining kind, so \(\mu (B)=b\). For each original \(A_i\), the finite-union supremum bounds \(\mu (F_n\cup A_i)\le b\), and continuity from below gives \(\mu (B\cup A_i)\le b=\mu (B)\). Thus \(A_i\setminus B\) is null. If \([C]\) is another upper bound, every one of the countably many constituent sets in \(B\) is contained in \(C\) modulo a null set; hence \(B\setminus C\) is null. This proves that \([B]\) is the supremum. Infima follow by complements.

In a complete Boolean algebra finite meets distribute over arbitrary joins: if \(v=\bigvee_i(a\wedge b_i)\), then every \(b_i\le \neg a\vee v\); taking joins and intersecting with \(a\) gives \(a\wedge\bigvee_i b_i\le v\). The reverse inequality is immediate. Thus \(L\) is a frame.

For a nonzero \([A]\), the function \(F(t)=\mu (A\cap [0,t])\) is continuous because \(|F(t)-F(u)|\le |t-u|\). It takes the values \(0\) and \(\mu(A)>0\) at the endpoints. There is therefore \(t\) with \(F(t)=\mu (A)/2\). The class of \(A\cap [0,t]\) is strictly between \(0\) and \([A]\), so no nonzero class is an atom. \(\square\)

\textbf{Theorem M2 --- exact frame points of a complete Boolean algebra.} If \(B\) is any complete Boolean algebra, its frame points correspond bijectively to its atoms:

\[
\operatorname{Hom}_{\mathrm{Frm}}(B,\mathbb B)
\cong\operatorname{At}(B),\qquad
a\longmapsto\epsilon_a(b)=\mathbf1_{\{a\le b\}}.
\tag{M1}
\]

\textbf{Proof.} An atom below \(\bigvee_i b_i\) lies below some \(b_i\): otherwise every \(a\wedge b_i=0\), while frame distributivity gives \(a=a\wedge\bigvee_i b_i=\bigvee_i(a\wedge b_i)=0\), a contradiction. Thus \(\epsilon _a\) preserves arbitrary joins. It preserves finite meets, top and bottom directly.

Conversely, a frame map \(\epsilon :B\to \mathbb B\) into the two-element Boolean algebra preserves complements, since complements are characterized by their finite meet \(0\) and join \(1\). By De Morgan it therefore also preserves arbitrary meets. Put \(a=\bigwedge\{b:\epsilon(b)=1\}\). Its image is \(1\), so \(a\ne 0\). If \(0<c\le a\) and \(\epsilon (c)=0\), then \(\epsilon (\neg c)=1\), forcing \(a\le \neg c\) and hence \(c=0\), a contradiction. Consequently \(\epsilon (c)=1\), which by definition of \(a\) implies \(a\le c\). Thus \(a\) is an atom. Every element of the filter contains \(a\); conversely if \(a\le b\), then \(\epsilon (b)=1\). This proves the inverse formulas and their uniqueness. \(\square\)

In particular the measure algebra of Lemma M1 has \textbf{no frame points}.

It nevertheless has ordinary lattice characters. Here is an explicit witnessed failure of the required continuity. Let \(D_n=[0,2^{-n}]\) as original measurable subsets of \([0,1]\); their classes are nonzero and any finite intersection is the smallest selected \(D_n\). The filter generated by these classes is proper. Order proper filters containing it by inclusion; the union of every chain is proper, so Zorn's lemma gives a maximal proper filter \(F\). In a Boolean algebra, maximal proper filters contain exactly one of \(b,\neg b\): if neither lay in \(F\), adjoining \(b\) would produce \(0\), so some \(f\in F\) satisfies \(f\wedge b=0\) and hence \(f\le \neg b\), which already puts \(\neg b\in F\). Its characteristic function is therefore a bounded-lattice homomorphism \(\epsilon_F:L\to\mathbb B\).

Every \(\epsilon _F([D_n])=1\), while every \(\epsilon _F([D_n^c])=0\). But

\[
\bigvee_{n\ge1}[D_n^c]=[[0,1]\setminus\{0\}]=1_L,
\quad\epsilon_F(1_L)=1\ne\bigvee_{n\ge1}0=0.
\tag{M2}
\]

Thus (B2) gives an actual surjective semiring branch \(\beta_{\epsilon_F}:G_L(R)\to G(R)\) preserving amplitudes, even though the frame has no point. Formula (M2) gives its exact failed zero-fibre join equation. The point \(\widehat\epsilon_F\) of \(\operatorname{Idl}(L)\) does not descend through \(q\). More explicitly, set

\[
J=\bigvee_{n\ge1}^{\operatorname{Idl}(L)}\downarrow[D_n^c].
\]

Then \(q(J)=1_L=q(L)\), but \(\widehat\epsilon_F(J)=0\) and \(\widehat\epsilon_F(L)=1\). This is an explicit pair in the kernel congruence of \(q\) distinguished by the upstairs point; it proves the obstruction to descent at an exact source/target pair.

In fact \textbf{every} ordinary branch on this measure algebra fails already countable-join preservation. For its ultrafilter \(F\), start at \(a_0=1\). Split a representative of \(a_n\) into two equal-measure pieces by the continuous cumulative-measure construction in Lemma M1. Exactly one of their classes belongs to \(F\); call it \(a_{n+1}\). The resulting decreasing sequence satisfies \(a_n\in F\) and \(\mu(a_n)=2^{-n}\). Its infimum is \(0\) by continuity of measure. Consequently \(b_n=\neg a_n\) increases to \(1\) and each \(\epsilon _F(b_n)=0\), again contradicting countable-join preservation. The ideal \(I=\bigcup_n\downarrow b_n\) then satisfies \(q(I)=q(L)=1\), \(\widehat\epsilon_F(I)=0\), and \(\widehat\epsilon_F(L)=1\). This gives a kernel-congruence witness for each individual algebraic branch, not merely for the displayed dyadic choice.

\subsection{Atomic observation and the exact support data it removes}\label{historical-support-atomic-observation-and-the-exact-support-data-it-removes}

For a complete Boolean algebra \(B\), put

\[
b_{\mathrm{at}}=\bigvee_{a\in\operatorname{At}(B)}a,\qquad
b_{\mathrm{na}}=\neg b_{\mathrm{at}}.
\]

The frame-point evaluation map is

\[
\operatorname{ev}:B\to\mathcal P(\operatorname{At}(B)),\qquad
x\mapsto\{a:a\le x\}.
\tag{A1}
\]

\textbf{Theorem A1 --- exact atomic quotient.} This is a surjective complete Boolean-algebra homomorphism. Its kernel congruence is

\[
\operatorname{ev}(x)=\operatorname{ev}(y)
\iff x\wedge b_{\mathrm{at}}=y\wedge b_{\mathrm{at}},
\tag{A2}
\]

and its inverse image of the empty set is the full principal ideal \(\downarrow b_{\mathrm{na}}\). The complete decomposition retaining that ideal is

\[
B\xrightarrow{\sim}(\downarrow b_{\mathrm{at}})\times(\downarrow b_{\mathrm{na}}),
\quad x\mapsto(x\wedge b_{\mathrm{at}},x\wedge b_{\mathrm{na}}),
\tag{A3}
\]

with inverse \((u,v)\mapsto u\vee v\) and respective local units \(b_{\mathrm{at}},b_{\mathrm{na}}\).

\textbf{Proof.} Each coordinate of (A1) is the frame point in (M1), so it preserves arbitrary joins and finite meets. Complements are preserved because an atom lies below exactly one of \(x,\neg x\). For any subset \(S\) of atoms, the element \(\bigvee_{a\in S}a\) maps exactly to \(S\), by pairwise disjointness of distinct atoms and the arbitrary-join argument in (M1). Thus evaluation is surjective. Distributivity gives \(x\wedge b_{\mathrm{at}}=\bigvee_{a\in\operatorname{At}(B)}(x\wedge a)=\bigvee_{\substack{a\in\operatorname{At}(B)\\a\le x}}a\), which proves (A2) and the empty-fibre claim. The two complementary components in (A3) reconstruct \(x\) by distributivity, and disjointness of \(b_{\mathrm{at}},b_{\mathrm{na}}\) proves both inverse identities and preservation of complete Boolean operations. \(\square\)

The complement on either local Boolean factor \(\downarrow b\) in (A3) is explicitly \(u\mapsto b\wedge \neg u\); its local top is \(b\), not the global top of \(B\) unless \(b=1\).

This supplies an exact relation to the source's adelic measure-algebra support. In a finite measure-algebra realization, frame-point evaluation detects precisely its measure atoms;(A3) retains the entire nonatomic component rather than treating lack of points as lack of support. For a discrete countable measure with positive measure at each singleton, the singleton classes are the atoms and \(b_{\mathrm{at}}=1\); evaluation is an isomorphism to the power set of the original places. For the Lebesgue example, \(b_{\mathrm{at}}=0\), \(b_{\mathrm{na}}=1\), and every frame-point evaluation is absent while the entire nontrivial Boolean carrier persists.

When applying \(G_{\operatorname{ev}}(R)\), retain the exceptional case: if there are no atoms, \(\mathcal P(\operatorname{At}(B))\) is the one-element lattice, so the exact carrier \(G_{\mathcal P(\varnothing)}(R)\) is isomorphic to \(R\), and this quotient identifies \(\tau _B\) and \(e_B\) with ordinary zero. It is the explicit ring-reflection map on this observation, not evidence that \(\tau _B=e_B\) in the original carrier. The split comparison (I7) has nontrivial lattices on both sides and preserves them as distinct; it remains available when the point-observation quotient collapses them.

Concretely, \(\Psi:G_{\{*\}}(R)\to R\), \((r,*)\mapsto r\), is bijective and preserves both operations because the unique support coordinate is fixed. In the atomless case its composition is the exact identity

\[
\Psi\circ G_{\mathrm{ev}}(R)=p_B:G_B(R)\to R,
\quad p_B(r,\lambda)=r.
\tag{A4}
\]

\subsection{The exact operator-algebraic branch map}\label{historical-support-the-exact-operator-algebraic-branch-map}

The core manuscript's \(L^\infty\) realization permits an additional precise bridge to state spaces. Work with the same finite probability measure space, complete measure algebra \(B\), and the complex algebra \(M=L^\infty(\Omega,\mu;\mathbb C)\) equipped with essential-supremum norm and complex conjugation. The original measure and every atomic mass are retained.

\textbf{Theorem O1 --- all Boolean branches extend exactly to multiplicative states.} There is a natural bijection

\[
\operatorname{Hom}_{\mathrm{DLat}_{0,1}}(B,\mathbb B)
\cong\{\chi:M\to\mathbb C:\chi\text{ is a unital complex-linear }*\text{-homomorphism}\}.
\tag{O1}
\]

The correspondence is specified on every projection by

\[
\chi_\epsilon(\mathbf1_A)=\epsilon([A]),
\tag{O2}
\]

and on arbitrary \(f\in M\) by the uniform extension of this finite formula:

\[
\chi_\epsilon\left(\sum_{j=1}^n c_j\mathbf1_{A_j}\right)
=\sum_{j=1}^n c_j\epsilon([A_j]),
\tag{O3}
\]

where the \(A_j\) form the specified measurable partition of \(\Omega\) modulo null sets. Every such \(\chi\) is positive and has norm \(1\).

\textbf{Proof.} Exactly one member of a finite partition has \(\epsilon\)-value \(1\): finite-join preservation gives at least one, and pairwise disjointness with meet preservation gives at most one. Thus (O3) selects one of the original values \(c_j\). Refining two partitions by all intersections shows the formula is independent of the simple-function presentation. The unique selected intersection gives additivity and multiplicativity for two simple functions, while conjugating their values gives star preservation. The selected value has modulus at most the essential supremum of the simple function, so the map is contractive.

Every bounded measurable complex function has simple functions tending to it in essential-supremum norm: partition a square containing its essential range into finitely many squares of diameter at most \(2^{-n}\) and use one original value in each nonempty square. The resulting evaluations are Cauchy by contractivity. Their limit defines a unique continuous extension independent of the approximating sequence. Addition, multiplication and conjugation commute with these uniform limits, proving the required homomorphism identities. Positivity follows from \(f=g^*g\) with \(g=\sqrt f\) for every nonnegative measurable \(f\); hence \(\chi(f)=|\chi(g)|^2\ge0\). Contractivity and \(\chi (1)=1\) prove norm \(1\).

Conversely a unital star homomorphism sends projections to real idempotents \(0\) or \(1\). Since \(\mathbf1_{A\cap B}=\mathbf1_A\mathbf1_B\) and \(\mathbf1_{A\cup B}=\mathbf1_A+\mathbf1_B-\mathbf1_{A\cap B}\), its projection values form the required Boolean lattice map. It is contractive because \(\|f\|_\infty^2\mathbf1-f^*f\) is the square of a bounded measurable function, so positivity gives \(|\chi(f)|^2\le\|f\|_\infty^2\). It agrees with (O3) on simple functions, and continuity forces equality on all of \(M\). This proves both inverse constructions. \(\square\)

\textbf{Theorem O2 --- normal branches, exact atomic masses, and failure of normality.} For the correspondence (O1), the following are equivalent:

\begin{enumerate}
\def\labelenumi{\arabic{enumi}.}
\tightlist
\item
  \(\chi\) preserves suprema of bounded increasing nets of nonnegative elements of \(M\) (normality).
\item
  \(\epsilon\) is a frame point of the complete measure algebra.
\item
  There is a measure atom \(a\) such that
  \[
  \chi(f)=\frac{1}{\mu(a)}\int_a f\,d\mu\qquad(f\in M).
  \tag{O4}
  \]
\end{enumerate}

In particular the atomless Lebesgue example has multiplicative states for all of its ordinary branches, and none is normal; its exact countable failure is exhibited by (M2).

All bounded increasing nets used here have a supremum in the specified \(M\). An explicit proof for nonnegative nets retains the finite measure: put \(b=\sup_i\int_\Omega f_i\,d\mu\). Choose indices with integrals tending to \(b\), and use directedness to replace them by an increasing dominating sequence. Choose representatives outside one common null set for the countably many sequence comparisons, so this selected sequence is pointwise increasing there. Its pointwise supremum \(g\) has integral \(b\). For every fixed original index \(i\), each finite maximum with a sequence member lies below a net member, so \(\int_\Omega\max(g,f_i)\,d\mu\le b=\int_\Omega g\,d\mu\) by monotone convergence. Thus \(f_i\le g\) almost everywhere. Every upper bound of the original net bounds the selected sequence and hence \(g\), proving that \(g\) is the required least upper bound.

\textbf{Proof.} The supremum in \(M\) of the finite-union projections from an arbitrary family \(([A_i])\) is the projection \(\mathbf1_{\bigvee_i[A_i]}\): Lemma M1's countable essential union is an upper bound, and its increasing finite selected unions show it is the least upper bound also among arbitrary measurable nonnegative functions. Normality applied to this increasing net therefore gives preservation of its full projection join. Together with (O2), this is precisely frame continuity of \(\epsilon\), proving \(1\Rightarrow2\). By Theorem M2, a frame point corresponds to a unique atom \(a\).

Every bounded measurable \(f\) is essentially constant on \(a\). To see this, use the finite measurable partitions of its essential range from the proof of Theorem O1. Each partition has exactly one cell whose intersection with \(a\) is nonzero, and by atomhood that intersection equals \(a\) modulo null sets. Its diameter tends to \(0\), so the simple approximants restrict to constants \(c_n\) on \(a\), with \(|c_n-c_m|\) bounded by their uniform difference. Their limit \(c\) is the almost-everywhere value of \(f\) on \(a\). Formula (O3) selects exactly these constants and gives \(\chi (f)=c\); integration over the original atom gives precisely (O4), with the factor \(1/\mu (a)\) retained. Thus \(2\Rightarrow3\).

Finally let a bounded increasing net \(f_i\ge 0\) have supremum \(f\) in \(M\), and let \(c_i\) be its respective constants on \(a\). These constants increase; put \(c=\sup_i c_i\). Every \(f_i\) is bounded above by the function equal to \(c\) on \(a\) and to \(f\) off \(a\); hence the least-upper-bound property of \(f\) forces its constant on \(a\) to be at most \(c\). Conversely each \(f_i\le f\) forces \(c_i\le f|_a\), and taking their supremum gives the reverse inequality. Thus \(\chi(f)=\sup_i\chi(f_i)\). This proves \(3\Rightarrow1\), including every original coefficient and atomic mass. \(\square\)

Equations (B1),(O1),(O2) form the promised typed connection between algebraic Split-Zero branches and operator-algebraic states. They identify exactly which branches can be normal observations of the measured support.
